\documentclass[10pt,a4paper]{amsart}
\usepackage[margin=19mm]{geometry}
\usepackage{amsmath,amssymb,mathtools}
\usepackage[scr=rsfs,cal=euler]{mathalfa}
\usepackage{xfrac}
\usepackage[unicode,hidelinks,bookmarks=false]{hyperref}
\newcommand{\ie}{i.e.\ }
\newcommand{\cf}{cf.\ }
\newcommand{\resp}{respectively\ }
\newcommand{\Subsection}{\subsection}
\providecommand{\nobreakdash}{\nobreak\hbox{-}}
\setlength{\emergencystretch}{2em}
\multlinegap0pt
\usepackage{bm}
\usepackage{bbm}
\usepackage{dsfont}
\usepackage{xspace}
\usepackage{cancel}
\usepackage{mathtools}
\usepackage{amsthm}
\makeatletter
\@ifundefined{definition}{\newtheorem{definition}{Definition}}{}
\@ifundefined{proposition}{\newtheorem{proposition}[definition]{Proposition}}{}
\makeatother
\newcommand{\calb}{\mathcal{B}}
\title[A characterization of locally ordered ternary relations in t]{A characterization of locally ordered ternary relations in terms of digraphs}
\author{Guillermo Gamboa Quintero, Mart\'in Matamala, Juan Pablo Pe\~na}
\date{Source: arXiv:2509.11956v1 (CC BY 4.0)}
\begin{document}
\maketitle

\noindent\emph{Source and licence.} A characterization of locally ordered ternary relations in terms of digraphs. Version arXiv:2509.11956v1 (CC BY 4.0), distributed under the Creative Commons Attribution 4.0 International licence, as recorded in the arXiv licence metadata for this exact version and shown on the abstract page https://arxiv.org/abs/2509.11956v1. Full rights notice: licenses/arxiv-2509.11956-CC-BY-4.0.txt.\par\smallskip

\noindent\textit{Editorial scope.} Counted unit: the Proposition whose source label is \texttt{p:tools} in the article (source0.tex lines 186--207, source label \texttt{p:tools}), delivered together with the article's own complete original proof of it (source0.tex lines 210--241, one \texttt{proof} environment of 32 lines). No mathematics is written, repaired or re-derived: statement and proof are the article's own text line for line. Required context retained uncounted: none. Every definition, display and label of those lines is retained unchanged. Omitted from this package and not redistributed: 1--185, 208--209, 242--350. Nothing inside a retained excerpt is omitted or altered: every retained body is delivered exactly as the article's lines are written, so no omission receipt of any kind accompanies this package. Citations: the retained excerpts carry no citation command at all, so no bibliography entry of the article is transcribed and no reference is redistributed. Reference adaptation: none was needed and none was made; every reference inside the delivered unit resolves inside it, so no unresolved reference or citation is produced. Printed slips kept literal: no printed word, symbol, hypothesis or number of the counted statement or of its proof was altered. Layout and class: the article's own class is \texttt{article} with packages including tikz, pstricks, epsfig, pst-grad, pst-plot, grffile, etoolbox, babel, geometry, amsmath, amsfonts, amssymb, amsthm, graphicx; the wrapper is the lane's pinned \texttt{amsart} editorial wrapper at 10pt on A4, which restates the theorem-like environments the retained text needs and adds the article's own macro definitions that the retained text needs (\texttt{\textbackslash calb}), with extra wrapper packages (bm, bbm, dsfont, xspace, cancel, mathtools). The delivered numbering is the unit's own. Assets: no retained span contains a figure environment, a graphics-inclusion command, a PDF-page inclusion, a table or a TikZ picture; no image file is copied into this package, the package declares no asset dependency and no image file is redistributed with it. Licence: version arXiv:2509.11956v1 of this article is distributed under the Creative Commons Attribution 4.0 International licence, as recorded in the arXiv licence metadata for this exact version and shown on the abstract page https://arxiv.org/abs/2509.11956v1. A full rights notice is delivered at licenses/arxiv-2509.11956-CC-BY-4.0.txt. Characters: every retained excerpt is pure ASCII, so the delivered engine, XeLaTeX with its default Latin Modern text font, reports no missing character.
\begin{proposition}
\label{p:tools}
	Let $\calb$ be a ternary relation on $V$ satisfying axioms \textbf{B},\textbf{C},\textbf{D} and \textbf{F}. Then, we have the following properties:
    
	\begin{enumerate}
		\item[\textbf{L}\,] $: \forall a,b,c \in V, abc\in\calb\implies \calb\cap \{a,b,c\}^3\in \{ \{abc,cba\}, \{abc,bca,cab\} \}.$ 
        
		\item[\textbf{G}\,]$: \forall a,b,c,d \in V, abc\notin\calb\implies bac \in\calb\lor acb\in\calb.$
	        
	    \item[\textbf{C'}]$: \forall a,b,c \in V, abc\in\calb\implies bac\notin\calb$
	\end{enumerate}
	
	Moreover, if \textbf{2} holds, then we also have the following properties:
    
	\begin{enumerate}
		\item[\textbf{3}\ ]$: \forall a,b,c,x \in V, abc\in\calb\land bxc \in\calb\implies axc\in\calb$.
	    \item[\textbf{2'}]$: \forall a,b,c,x \in V, abc\in\calb\land axb \in\calb\implies axc\in\calb$.
	    \item[\textbf{3'}]$: \forall a,b,c,x \in V, abc\in\calb\land axb \in\calb\implies xbc\in\calb$.
            
            
	\end{enumerate}
\end{proposition}

\begin{proof}
	To ease the presentation, along the proof we do not mention explicitely the use of axiom \textbf{D}. Let us assumme that \textbf{B}, \textbf{C}, \textbf{D} and \textbf{F}. 

	\begin{enumerate}
		\item[\textbf{L}\,] As $abc\in\calb$, by \textbf{F} we have that $cba\in\calb\lor bca\in\calb$. Assumme $cba\in\calb$.   By \textbf{C} we get that $cab\notin\calb\land acb\notin\calb$. Now, we have to prove that $ bca\notin\calb\land bac\notin\calb$. Suppose that $bca\in\calb$, by \textbf{F} we get $cab\in\calb\lor acb\in\calb$, which is a contradiction. Thus, $bac\in\calb$.   By \textbf{F}, $acb\in\calb\lor cab\in\calb$, which is also a contradiction. Then  $ \calb\cap \{a,b,c\}^3=\{abc,cba\}$.
	
		Now, assume $bca\in\calb$.   By \textbf{C}, we have that $ bac\notin\calb\land acb\notin\calb$.   By \textbf{F}, we have that $acb\in\calb\lor cab\in\calb$, since $acb\notin\calb$ then $cab\in\calb$ and $ cba\notin\calb$. Thus, $ \calb\cap\{a,b,c\}^3=\{abc,bca,cab\}$.
        
		\item[\textbf{G}\,] Assume that $ abc\notin V \land bac\notin V\land acb\notin\calb$.   By $\textbf{B}$ we get that $bca\in\calb\lor cab\in\calb\lor cba\in\calb$ and by \textbf{L} we know that $| \calb\cap \{a,b,c\}^3|=3$. Thus, $ \calb\cap \{a,b,c\}^3=\{bca,cab,cba\}$ which is a contradiction, as $cab\in\calb\implies cba\notin\calb$, by \textbf{C}.
	 
		\item[\textbf{C'}\,] Since $abc\in\calb$, by \textbf{L}  we get that $ \calb\cap \{a,b,c\}^3$ is equal to  $\{abc,cba\}$ or to $\{abc,bca,cab\}$. Either case, $bac\notin\calb$.
        	\end{enumerate}

	Now, assume that \textbf{2} holds.
	\begin{enumerate}
		\item[\textbf{3}\,] We want to prove $axc\in\calb$.   By \textbf{2}, $abc\in\calb\land bxc\in\calb\Rightarrow abx\in\calb$.  By \textbf{L} we know that the set of valid triples of the set $\{b,c,x\}$ is either $\{bxc,xcb,cbx\}$ or $\{bxc,cxb\}$. %\textcolor{blue}{done}\textcolor{red}{We need to define what valid means, if we are using it.}
	
		In the first case, by \textbf{2}, we have that $acx\in\calb\land cbx\in\calb \Rightarrow acb\in\calb$ which contradicts $abc\in\calb$. Hence, $acx\notin\calb$ and, by \textbf{G} we get that $cax\in\calb\lor axc\in\calb$. If $cax\in\calb$, then, by \textbf{2} applied to $cax\in\calb\land abx\in\calb$, we get that $cab\in\calb$. Again, by \textbf{2} and $xcb\in\calb\land cab\in\calb$, we get $xca\in\calb$. From \textbf{L} applied to $\{a,c,x\}$ we conclude $axc\in\calb$. 
	
		In the second case, we first prove that $ bxa\notin\calb$. Otherwise, from \textbf{L}  we  also get $xab\in\calb$. By \textbf{2} applied to $cxb\in\calb\land xab\in\calb$ we get $cxa\in\calb$. Again by \textbf{2} applied to $bca\in\calb\land cxa\in\calb$ and to $cba\in\calb\land bxa\in\calb$ we get $bcx\in\calb$ and $cbx\in\calb$, respectively. Thus, we get  $bca\notin\calb\land cba\notin\calb$ which contradicts \textbf{L} applied to $\{a,b,c\}$.
	
		Since $bxa\notin\calb$, by \textbf{L} applied to $\{a,b,x\}$ we get $xba\in\calb$.   By \textbf{2}, we have that $bxc\in\calb\land xac\in\calb\Rightarrow bxa\in\calb$ which is a contradiction with  \textbf{L} applied to $\{a,b,x\}$. Thus, we have $xac\notin\calb$ and by \textbf{C'} we obtain $axc\in\calb\lor xca\in\calb$. When $xca\in\calb$ we can apply \textbf{2} to $xca\in\calb\land cba\in\calb$ to get $xcb\in  \calb$, which is contradiction. Then, we have $xca\notin  \calb$ and thus the desired conclusion $axc\in  \calb$ holds.
	
		\item[\textbf{2'}\,] Let $abc\in  \calb\land axb\in  \calb$. Assume that $axc \notin  \calb$.   By \textbf{G} we have that $xac\in  \calb$ or $acx\in  \calb$.   By \textbf{3}, $xac\in  \calb\land abc\in  \calb\implies xab\in  \calb$ which contradicts $axb\in  \calb$, by \textbf{C'}.
	
		Thus, we have that $xac\notin\calb\land acx\in  \calb$.   By \textbf{L} we get that $xca\in  \calb$.   By \textbf{2},  $xca\in  \calb\land cba\in\calb\implies xba\in  \calb$. Again by \textbf{L} we get that $bxa\notin  \calb$ and thus $bax\in  \calb$.   By \textbf{3} we get $bax\in  \calb\land acx\in  \calb\implies bac\in  \calb$ which is a contradiction with $abc\in  \calb$, by \textbf{C'}.
	
		\item[\textbf{3'}]  We prove \textbf{3'}. Assume $ xbc \notin  \calb$. Then \textbf{G} implies $bxc\in  \calb\lor xcb\in  \calb$. In the first case, by \textbf{3}, $abc\in  \calb\land bxc\in  \calb\implies abx\in  \calb$, which is a contradiction with $axb\in  \calb$. Thus, $xcb\in  \calb$.   By \textbf{2} we have that $axb\in  \calb\land xcb\in  \calb\implies acb\in  \calb$ which is a contradiction with $abc\in  \calb$.
        
        
	\end{enumerate}
 \end{proof}

\end{document}
